%%%%%%%%%%%%%%%%%%%%%%%%%%%%%%%%%%%%%%%%%%%%%%%%%%%%%%%%%%%%%%%%%%%%%%%%%%%%%%%%
%% 

\chapter{ASP-Logik}
In diesem Kapitel wird genauer auf die ASP-Logik und die Generierung der Dienstpläne eingegagen.
\section{Fakten-Generierung}
Die Fakten-Generierung stellt die Verbindung zwischen der relationalen Datenbank und der ASP-Wissensbasis dar. Damit Clingo einen Dienstplan berechnen kann, müssen die benötigten Informationen zunächst in Form von ASP-Fakten bereitgestellt werden. Dafür werden die aus der Datenbank geladenen ORM-Objekte in der Funktion \texttt{\_build\_base\_facts} in entsprechende Fakten umgewandelt.

Für jeden Benutzer wird dabei ein \texttt{user}-Fakt erzeugt. Zusätzlich wird für jede vorhandene Kompetenz ein \texttt{has\_skill}-Fakt erstellt. Für Stationen werden ebenfalls Fakten generiert, darunter \texttt{station}, \texttt{max\_assignments} zur maximalen Anzahl an Zuweisungen sowie \texttt{requires} für benötigte Kompetenzen.

Die zu planenden Arbeitstage werden als \texttt{day}-Fakten an Clingo übergeben. Wochenenden werden dabei bereits vorher durch die Funktion \texttt{get\_weekdays} herausgefiltert. Zusätzlich werden benutzerspezifische Einschränkungen als Fakten festgelegt. Dazu zählen beispielsweise maximale oder minimale Arbeitstage pro Monat (\texttt{max\_days}, \texttt{min\_days}), fixe Arbeitstage (\texttt{fixed\_day}) oder gesperrte Tage (\texttt{blocked\_day}).

Im Folgenden wird beispielhaft ein Ausschnitt der generierten Basisfakten für ein \textit{Scheduling} gezeigt. 

\begin{lstlisting}[caption={Beispiel generierter Basisfakten},label={lst}]
	user(1597).
	has_skill(1597, 23).
	has_skill(1597, 24).
	
	user(1598).
	has_skill(1598, 22).
	has_skill(1598, 25).
	
	station(865).
	max_assignments(865, 3).
	requires(865, 22).
	requires(865, 23).
	
	station(866).
	max_assignments(866, 3).
	requires(866, 24).
	
	day(1).
	day(4).
	day(5).
	
	exact_days(1597, 21).
	exact_days(1598, 21).
\end{lstlisting}
In diesem Beispiel werden zwei Benutzer, deren Kompetenzen, Stationen mit den benötigten Kompetenzen sowie Arbeitstage und Constraints als Fakten definiert. Auf Basis dieser Informationen kann Clingo anschließend gültige Dienstpläne berechnen.

\section{Regelwerk}
%First Draft%
Das Regelwerk für die initiale Planerstellung ist in der Datei \texttt{rules.lp} 
definiert. Die Grundlage bildet die Bestimmung der \textit{Eligibility}: Ein 
Benutzer gilt als qualifiziert für eine Station (\texttt{eligible/2}), wenn er 
alle erforderlichen Fähigkeiten besitzt, also kein \texttt{requires/2}-Fakt 
existiert, dem kein entsprechendes \texttt{has\_skill/2}-Fakt gegenübersteht. 
Daraus wird \texttt{eligible\_day/3} abgeleitet, welches zusätzlich gesperrte 
Tage ausschließt. Der \textit{Entscheidungsraum} wird durch eine \textit{Choice 
	Rule} definiert, die Clingo erlaubt, Benutzern Stationen an bestimmten Tagen 
zuzuweisen, wobei pro Station und Tag maximal \texttt{Max} Benutzer zugewiesen 
werden dürfen. Die \textit{harten Constraints} erzwingen, dass kein Benutzer an 
zwei Stationen gleichzeitig eingeteilt wird, dass gesperrte Tage nicht belegt 
werden, und dass benutzerspezifische Einschränkungen wie \texttt{max\_days/2}, 
\texttt{min\_days/2} sowie \texttt{exact\_days/2} eingehalten werden. Fixe 
Arbeitstage werden über \texttt{fixed\_day/2} als harte Bedingung erzwungen, 
sodass bestimmte Benutzer an vordefinierten Tagen zwingend eingeteilt sind. 
Die \textit{Optimierung} minimiert unbesetzte Stationen über eine 
\texttt{\#minimize}-Direktive, die Clingo anweist, die Anzahl der Tage, an 
denen eine Station ohne Besetzung bleibt, so gering wie möglich zu halten. 
Die Regelstruktur für das Rescheduling folgt einem ähnlichen Prinzip, wird 
jedoch in Abschnitt~\ref{sec:rescheduling} gesondert behandelt.
\section{Zweiphasen-Lösungsanstz}

\section{Rescheduling-Logik}

\section{Timeout NP-hartes Problem}



%% 
%%%%%%%%%%%%%%%%%%%%%%%%%%%%%%%%%%%%%%%%%%%%%%%%%%%%%%%%%%%%%%%%%%%%%%%%%%%%%%%%
